% fig_vl18_statistik.tex — Was das Vorzeichen in P_- kostet: die
% gemeinsame Aufenthaltsdichte zweier identischer Teilchen im Kasten,
% symmetrisiert gegen antisymmetrisiert.
%
% WHY THIS FIGURE.  The lecture proves that H^{\otimes N} splits into
% H_+ and H_-, and that the sign of U_pi cancels in rho.  What it never
% shows is what the sign DOES.  It does something drastic, and it is a
% statement about a two-dimensional shape, which is exactly what prose
% handles badly: bosons pile up on the diagonal x_1 = x_2, fermions are
% expelled from it, and the two densities are otherwise the same lumps
% in mirrored places.
%
% GEOMETRY FROM THE FORMULA.  Nothing is drawn by eye.  The two
% particles sit in the box of Vorlesung 11, orbitals n = 1 and n = 2,
%
%       psi_n(x) = sqrt(2/L) sin(n pi x / L),
%
% and the two-particle vectors are the ones the projectors P_\pm of this
% lecture produce out of psi_1 (x) psi_2 (x):
%
%   Psi_\pm(x_1,x_2) = [psi_1(x_1)psi_2(x_2) \pm psi_2(x_1)psi_1(x_2)]/sqrt2,
%
% so with L = 1
%
%   |Psi_\pm|^2 = 2 [ sin(pi x_1) sin(2pi x_2) \pm sin(2pi x_1) sin(pi x_2) ]^2.
%
% Expanding the square splits it into a direct part and an exchange
% part; the direct part alone is the distinguishable case
%
%   P_u = (|Psi_+|^2 + |Psi_-|^2)/2,
%
% i.e. the grey curve of panel (c) is the exact arithmetic mean of the
% blue and the red one.  The exchange term is the whole difference.
%
% The exact numbers the plate is built on:
%   * both densities have exactly two maxima, of the SAME height
%         max |Psi_\pm|^2 = 128/27 = 4.7407... ,
%     attained at cos(pi x) = 1/sqrt3, i.e. at
%         x^* = arccos(1/sqrt3)/pi = 0.304087 ,   1 - x^* = 0.695913 ;
%     for Psi_+ they sit at (x^*,x^*) and (1-x^*,1-x^*) — ON the
%     diagonal — for Psi_- at (x^*,1-x^*) and (1-x^*,x^*) — OFF it.
%     Same two lumps, moved.  That is why one colour scale serves both
%     maps and both actually reach its top.
%   * the distinguishable density peaks at 64/27, exactly half.
% Everything is therefore plotted in units of 64/27 (factor 27/64 =
% 0.84375 in front of the brackets), which puts the three numbers that
% matter on the axis ticks: 2 (boson), 1 (distinguishable), 0 (Pauli).
%
% Panel (c) is the horizontal dashed cut x_2 = x^* through both maps —
% particle 2 pinned where the pair density is largest, particle 1 free.
% There
%       |Psi_+|^2 = 2, |Psi_-|^2 = 0  at coincidence x_1 = x^*,
%       |Psi_+|^2 = 0, |Psi_-|^2 = 2  at the mirror point 1 - x^*,
% with the distinguishable reference equal to 1 at both, and all three
% curves meeting at x_1 = 1/2.  So the antibunching is not a lack of
% probability: the weight fermions take off the diagonal reappears,
% undiminished, on the other side.
%
% DESIGN DECISIONS
%   * colour codes the SIGN of the exchange and nothing else: cgblue for
%     the symmetric state, birrenred for the antisymmetric one.  The
%     panel titles are set in those two colours, which is what makes the
%     legend in (c) unnecessary — the curves are identified by the
%     titles above the maps they belong to;
%   * inside a map colour is the density itself (the neutral tpcdensity
%     ramp), on the same scale in both, so the maps may be compared
%     lump for lump;
%   * dotted grey means coincidence x_1 = x_2 in every panel: the
%     diagonal of the maps and the vertical marker of (c) are the same
%     line.  Dashed grey means the cut, and it appears only where the
%     cut is taken;
%   * the distinguishable reference is grey and dashed because it is a
%     reference, not a third species;
%   * the maps are square, with equal x and y ranges, because the
%     exchange symmetry of |Psi|^2 about the diagonal is the point and a
%     stretched box would hide it.

\begin{center}
\begin{tikzpicture}

% ─────────────────────────────────────────────────────────────
%  Shared abbreviations.  x^* = arccos(1/sqrt3)/pi, and the values of
%  the two orbitals there: sin(pi x^*) = sqrt(2/3), sin(2 pi x^*) = 2sqrt2/3.
% ─────────────────────────────────────────────────────────────
\def\xstar{0.3040867}
\def\sxstar{0.8164966}   % sin(pi x^*)
\def\dxstar{0.9428090}   % sin(2 pi x^*)

% ─────────────────────────────────────────────────────────────
%  (a)  Bosons: the symmetric density, bunched on the diagonal
% ─────────────────────────────────────────────────────────────
\begin{axis}[
    name=bos,
    tpcplot,
    width=3.2cm, height=3.2cm, scale only axis,
    view={0}{90},
    colormap name=tpcdensity,
    point meta min=0, point meta max=2,
    grid=none,
    xmin=0, xmax=1, ymin=0, ymax=1,
    enlargelimits=false,
    xtick={0,1}, xticklabels={$0$,$L$},
    ytick={0,\xstar,1}, yticklabels={$0$,$x^{*}$,$L$},
    ztick=\empty,
    xlabel={$x_1$}, ylabel={$x_2$},
    every axis x label/.style={at={(0.5,-0.13)}, anchor=north, font=\small},
    every axis y label/.style={at={(-0.30,0.5)}, rotate=90,
                               anchor=south, font=\small},
    title={(a) \tpcifenglish{bosons}{Bosonen} $|\Psi_+|^2$},
    title style={text=cgblue, at={(0.5,1.04)}},
]
  \addplot3[surf, shader=interp, samples=61, domain=0:1, y domain=0:1]
    {0.84375*( sin(deg(pi*x))*sin(deg(2*pi*y))
             + sin(deg(2*pi*x))*sin(deg(pi*y)) )^2};

  % coincidence line x_1 = x_2 (dotted grey in every panel)
  \addplot3[classicallim, no marks] coordinates {(0,0,3) (1,1,3)};
  % the cut of panel (c)
  \addplot3[envelope, no marks] coordinates {(0,\xstar,3) (1,\xstar,3)};
\end{axis}

% ─────────────────────────────────────────────────────────────
%  (b)  Fermions: the same lumps, moved off the diagonal
% ─────────────────────────────────────────────────────────────
\begin{axis}[
    name=fer,
    at={(bos.south east)}, anchor=south west, xshift=9mm,
    tpcplot,
    width=3.2cm, height=3.2cm, scale only axis,
    view={0}{90},
    colormap name=tpcdensity,
    point meta min=0, point meta max=2,
    grid=none,
    xmin=0, xmax=1, ymin=0, ymax=1,
    enlargelimits=false,
    xtick={0,1}, xticklabels={$0$,$L$},
    ytick={0,\xstar,1}, yticklabels={,,},
    ztick=\empty,
    xlabel={$x_1$},
    every axis x label/.style={at={(0.5,-0.13)}, anchor=north, font=\small},
    title={(b) \tpcifenglish{fermions}{Fermionen} $|\Psi_-|^2$},
    title style={text=birrenred, at={(0.5,1.04)}},
    colorbar,
    colorbar style={width=0.17cm, ytick={0,1,2},
                    tick label style={font=\tiny, text=spacecadet},
                    axis line style={spacecadet, line width=0.4pt}},
]
  \addplot3[surf, shader=interp, samples=61, domain=0:1, y domain=0:1]
    {0.84375*( sin(deg(pi*x))*sin(deg(2*pi*y))
             - sin(deg(2*pi*x))*sin(deg(pi*y)) )^2};

  \addplot3[classicallim, no marks] coordinates {(0,0,3) (1,1,3)};
  \addplot3[envelope, no marks] coordinates {(0,\xstar,3) (1,\xstar,3)};

  \node[lbl, font=\scriptsize, text=spacecadet, anchor=center, rotate=45]
    at (axis cs:0.62,0.62,3) {\tpcifenglish{Pauli hole}{Pauli-Loch}};
\end{axis}

% ─────────────────────────────────────────────────────────────
%  (c)  The cut x_2 = x^*: 2 : 1 : 0 at coincidence
% ─────────────────────────────────────────────────────────────
\begin{axis}[
    at={(fer.south east)}, anchor=south west, xshift=21mm,
    tpcplot,
    width=3.9cm, height=3.2cm, scale only axis,
    xmin=0, xmax=1, ymin=0, ymax=2.62,
    xtick={0,\xstar,1}, xticklabels={$0$,$x^{*}$,$L$},
    ytick={0,1,2},
    grid=none,
    xlabel={$x_1$},
    ylabel={$|\Psi(x_1,x^{*})|^{2}$},
    every axis x label/.style={at={(0.5,-0.13)}, anchor=north, font=\small},
    every axis y label/.style={at={(-0.26,0.5)}, rotate=90,
                               anchor=south, font=\small},
    title={(c) \tpcifenglish{cut at}{Schnitt bei} $x_2=x^{*}$},
    domain=0:1, samples=201,
]
  % coincidence marker, same dotted grey as the diagonal of (a),(b)
  \draw[classicallim] (axis cs:\xstar,0) -- (axis cs:\xstar,2.62);

  % distinguishable reference = mean of the two
  \addplot[envelope]
    {0.84375*( sin(deg(pi*x))^2*\dxstar^2 + sin(deg(2*pi*x))^2*\sxstar^2 )};
  % symmetric and antisymmetric
  \addplot[wavefn, cgblue]
    {0.84375*( sin(deg(pi*x))*\dxstar + sin(deg(2*pi*x))*\sxstar )^2};
  \addplot[wavefn, birrenred]
    {0.84375*( sin(deg(pi*x))*\dxstar - sin(deg(2*pi*x))*\sxstar )^2};

  % the three values at coincidence
  \node[circle, fill=cgblue,    inner sep=1.3pt] at (axis cs:\xstar,2) {};
  \node[circle, fill=birrengray,inner sep=1.3pt] at (axis cs:\xstar,1) {};
  \node[circle, draw=birrenred, fill=white, line width=0.7pt,
        inner sep=1.1pt] at (axis cs:\xstar,0) {};

  \node[lbl, text=cgblue,    anchor=south east] at (axis cs:0.28,2.03) {$\Psi_+$};
  \node[lbl, text=birrenred, anchor=south west] at (axis cs:0.72,2.03) {$\Psi_-$};

  % inline key for the grey reference
  \draw[envelope] (axis cs:0.30,2.40) -- (axis cs:0.40,2.40);
  \node[lbl, text=birrengray, font=\scriptsize, anchor=west]
    at (axis cs:0.43,2.40) {\tpcifenglish{distinguishable}{unterscheidbar}};
\end{axis}

\end{tikzpicture}
\end{center}
